% ---
% Informações de dados para CAPA e FOLHA DE ROSTO
% Carregado por \input no PREÂMBULO do principal.tex: \data{} só pode ser usado
% no preâmbulo (dentro do documento o LaTeX acusa "Can be used only in preamble").
% ---
% Título definido pelo Enzo em 14/11/2026. Substitui o provisório, que repetia o tema do
% projeto FAPEMIG ("Estudo e desenvolvimento de estratégias de compressão de redes neurais
% para redução do consumo de energia em LLMs").
\titulo{Poda Guiada por Sensibilidade para Eficiência Energética aplicada ao modelo de linguagem GPT-2}
% Nome completo a pedido do Enzo em 16/11/2026. Antes: "Enzo Veloso".
\autor{Enzo Veloso Santos Costa}
% NORMALIZAÇÃO (14/11/2026): o manual pede o local como cidade (itens 4.1 e 4.2; o modelo
% da UFVJM traz "Diamantina"). Antes: "Diamantina/MG".
\local{Diamantina}
\data{2026}
\orientador{Prof. Dr. Marcelo Ferreira Rego}
%\coorientador{Equipe \abnTeX}
% NORMALIZAÇÃO (14/11/2026): universidade e curso separados, porque a capa do manual
% (item 4.1) os imprime em linhas e formatos próprios. Mesmos dizeres de antes.
\newcommand{\imprimiruniversidade}{UNIVERSIDADE FEDERAL DOS VALES DO JEQUITINHONHA E MUCURI}
\newcommand{\imprimircurso}{Bacharelado em Sistemas de Informação}
% Faculdade e departamento, impressos na capa entre a universidade e o curso (a pedido do
% Enzo em 16/11/2026). O item 4.1 do manual lista os elementos obrigatórios da capa e não
% veda outros; as duas linhas seguem a formatação do curso.
\newcommand{\imprimirfaculdade}{Faculdade de Ciências Exatas e Tecnológicas}
\newcommand{\imprimirdepartamento}{Departamento de Computação}
\instituicao{%
  \imprimiruniversidade
  \par
  \imprimircurso
  }
% Instituição de cada membro da banca na folha de aprovação (manual, item 4.5).
\newcommand{\imprimirinstituicaoorientador}{UFVJM}
% Avaliadores: os dois da UFVJM, informado pelo Enzo em 16/11/2026 (antes, "[TO DO] Instituição").
\newcommand{\imprimirinstituicaobancaum}{UFVJM}% Profa. Dra. Cinthya Rocha Tameirão
\newcommand{\imprimirinstituicaobancadois}{UFVJM}% Prof. Dr. Marcus Vinícius Carvalho Guelpeli
\tipotrabalho{Monografia (Graduação)}

% O preambulo deve conter o tipo do trabalho, o objetivo,
% o nome da instituição e a área de concentração
\preambulo{Monografia apresentada ao Curso de Sistemas de Informação do Departamento de Computação da Universidade Federal dos Vales do Jequitinhonha e Mucuri como requisito parcial para obtenção do título de Bacharel em Sistemas de Informação.}
% ---
